\begin{table}[t]
\centering
\caption{Bandwidth robustness: treatment effects $\Delta$ (pp, treatment $-$ baseline) within each bandwidth condition (primary model: Mistral Small Creative). Top row shows baseline join rates for reference.}
\label{tab:bandwidth}
\small
\begin{tabular}{lccc}
\toprule
 & BW=0.05 & BW=0.15 & BW=0.30 \\
\midrule
Baseline (level) & 0.054 & 0.409 & 0.061 \\
\midrule
\multicolumn{4}{l}{\textit{Treatment effect $\Delta$ (treatment $-$ baseline):}} \\
Stability & +0.7 & -1.4 & +0.9 \\
Upper cens. & +6.2 & -15.5 & +5.3 \\
Lower cens. & +10.1 & +3.7 & +9.6 \\
\bottomrule
\end{tabular}
\vspace{0.25em}
\parbox{\columnwidth}{\footnotesize\emph{Notes:} Primary model. All three bandwidth conditions use the same calibrated Mistral setup with $B = C = 1$ ($\theta^* = 0.50$, grid $[0.20, 0.80]$). $h = 0.15$: primary information design experiment. $h \in \{0.05, 0.30\}$: separate reruns with the same prompt template and \texttt{--load-calibrated} briefing parameters, but executed in dedicated batches (20 reps per cell under \texttt{output/bandwidth-005} and \texttt{output/bandwidth-030}). Their summaries also use an older output schema, omitting the briefing-diagnostic columns recorded in the main $h = 0.15$ run. Level comparisons across columns are therefore not meaningful; the low baselines reflect between-batch decoding variation rather than different payoff or grid settings. $\Delta$: change in mean join vs.\ within-condition baseline (pp).}
\end{table}
